\chapter{Programmazione di GPU e CUDA}

Una componente essenziale in tutti i computer è la CPU, che è nata essenzialmente per eseguire calcoli sequenziali nel più breve tempo possibile. Fin dalla loro creazione, l'idea era di \textbf{ottimizzare} quanto più possibile la \textbf{riduzione della latenza} tra un calcolo e l'altro, usando meccanismi come branch predictors, istruzioni atomiche di operazioni congiunte, etc... Un computer necessita di una CPU per poter funzionare e per garantire alcune funzioni base, tuttavia le sue capacità possono essere espanse attraverso altre componenti secondarie. Una di queste componenti è la \textbf{Graphics Processing Unit}, altresì detta \textbf{GPU}.

\begin{definition}{Graphics Processing Unit (GPU)}
    Una \textbf{Graphics Processing Unit} (\textbf{GPU}) è una \textbf{componente secondaria} di un computer, che venne inizialmente ideata per condurre compiti di grafica, ma che venne poi usata per \textbf{eseguire operazioni in parallelo} una volta scoperta la potenzialità di queste architetture
\end{definition}

La struttura di una GPU è molto diversa da quella di una CPU: mentre quest'ultima infatti è pensata per eseguire codice sequenziale, la prima è ideata per eseguire semplici operazioni in parallelo. In generale, si vuole \textbf{usare una CPU} quando si vuole avere una \textbf{latenza ridotta}, mentre invece si \textbf{usare una GPU} quando si vuole avere un \textbf{throughput maggiore}.

\section{CUDA}

In questo corso ci concentreremo sulle architetture delle GPU con architettura CUDA, create da NVIDIA. La società statunitense si è ormai affermata come leader mondiale nel campo del calcolo parallelo, sviluppando GPU e altre componenti che rappresentano lo stato dell'arte del campo del calcolo ad alte prestazioni; da qui il motivo del perché ci concentreremo sulle loro piattaforme.
\nl
Al giorno d'oggi però sono due i grandi competitors nel mercato delle GPU: NVIDIA e AMD. Sebbene NVIDIA sia la società più quotata, AMD offre comunque hardware performante, e rimane una valida alternativa a NVIDIA. AMD offre un'API proprietara chiamata HIP, che nella maggior parte delle istruzioni è simile a CUDA. Ci sono anche altre due API open source, chiamate OpenCL e OpenACC:
\begin{itemize}
    \item Open Computing Language (OpenCL) è uno standard di programmazione che punta allo sviluppare codice che possa essere eseguito su una varietà corposa di piattaforme, come CPU, GPU, DSP e altri tipi di processori. Sia NVIDIA che AMD supportano questo standard, ed è molto vicino alle API di CUDA;
    \item OpenACC è una specifica per un'API simile a OpenMP: l'idea è di usare delle pragma per parallelizzare i programmi sequenziali e farli eseguire su piattaforme che possano svolgere calcoli in parallelo, come le GPU.
\end{itemize}

Abbiamo spesso parlato di CUDA, ma di cosa parliamo quando usiamo questo termine? 

\begin{definition}{CUDA}
    \textbf{CUDA} è un acronimo che sta per \textbf{Compute Unified Device Architecture}, il quale denota l'\textbf{architettura} delle schede GPU di NVIDIA. CUDA offre due API diverse, che permettono di avere un controllo più granulare sulle piattaforme di NVIDIA
\end{definition}

CUDA va oltre il semplice linguaggio di programmazione: è una \textbf{piattaforma di calcolo} a tutti gli effetti. Le piattaforme CUDA vengono considerate come \textbf{co-processori} della CPU, non come sostituzioni. Ricordiamo che infatti, per funzionare, la GPU necessita della CPU.
\nl
Una GPU basata su CUDA è composta dalle seguenti componenti:
\begin{itemize}
    \item una \textbf{sezione ridotta di cache}, che puntano a massimizzare il throughput della memoria;
    \item una \textbf{sezione massiva di ALUs} (più lenti ed economici di quelli nelle CPU), che sono però collegati da una pipeline pensata apposta per svolgere compiti in parallelo;
    \item una memoria di tipo HBM (High Bandwidth Memory) chiamata VRAM (\textbf{Video RAM}), che è esclusiva alla GPU.
\end{itemize}

\begin{center}
    \includegraphics[scale = 0.27]{imgs/001.png}
\end{center}

L'architettura CUDA prevede la seguente struttura: ogni ALU è racchiuso in un core, chiamato \textbf{streaming processor} (\textbf{SP}); gli SP sono uniti in più blocchi, chiamati \textbf{streaming multiprocessors} (\textbf{SM}). Ogni streaming multiprocessor è controllato da una \textbf{stessa unità di controllo}: questo vuol dire che ogni streaming processor in uno stesso multiprocessor eseguirà la \textbf{stessa istruzione di Assembly} in un dato momento. Ogni SM ha anche una sua cache e una memoria \verb|texture|, che però non verrà esaminata.
\nl
Ci sono alcuni caveat quando si programmano le GPU: il primo è che la memoria della GPU e quella della CPU sono diverse, e quindi per poter usare un dato sulla GPU che risiede nella memoria della CPU bisognerà prima caricarlo sulla memoria della GPU; il secondo è che le GPU potebbero aderire a standard diversi, soprattutto quando si parla di numeri a virgola mobile. Questo si verifica generalmente quando bisogna arrotondare numeri particolari.

\subsection{Struttura di un programma CUDA}

CUDA organizza i suoi thread in una struttura a sei dimensioni: ogni thread è posizionato in un blocco che va da una a tre dimensioni, e ogni blocco è posizionato in una griglia che va anch'essa da una a tre dimensioni. Dunque ogni thread può sapere la sua posizione in questo spazio a sei dimensioni, chiamando opportune funzioni offerte dall'API di CUDA.
\nl
Ogni piattaforma CUDA ha una sua capacità di calcolo, che viene chiamata \textbf{CUDA compute capability}. La capacità viene rappresentata da un numero di versione chiamato \textbf{SM version}. Le tipologie di architetture compatibili con CUDA sono parecchie e variano da GPU a GPU: quelle più professionali e più mirate per compiti ad alta precisione avranno un valore di compute capability più alto rispetto alle GPU desktop che si possono trovare nei computer di fascia media-alta. 
\nl
Ora che sappiamo come i threads sono organizzati in un programma CUDA, possiamo iniziare a vedere come scrivere un programma. La struttura di un programma CUDA è la seguente:
\begin{itemize}
    \item [1)] \textbf{allocare la memoria} sulla piattaforma CUDA;
    \item [2)] \textbf{trasferire i dati} dalla DRAM alla VRAM della piattaforma;
    \item [3)] \textbf{eseguire un kernel CUDA}, ovverosia una funzione;
    \item [4)] una volta terminato il kernel, servirà \textbf{copiare i dati} processati dalla GPU; questi saranno disponibili sulla VRAM della piattaforma, e andranno copiati nella DRAM.
\end{itemize}

Per scrivere un programma di CUDA bisogna specificare una funzione, chiamata \textbf{kernel}, che andrà eseguita su tutti i threads. Chiaramente non è pensabile specificare, thread per thread, cosa ogni thread dovrà eseguire, qundi l'idea è quella di specificare come arrangiare i threads in blocchi e griglie e far eseguire a gruppi di threads la stessa funzione.
\nl
Di norma, la GPU ha già una sua struttura di blocchi e griglie, tuttavia noi abbiamo la possibilità di scegliere una nostra struttura che sia compatibile con l'architettura della nostra GPU. Ciò è possibile attraverso la specifica di blocchi / griglie custom, di tipo \verb|dim3|, che è un vettore di interi. Possiamo creare griglie e blocchi usando il seguente codice:

\begin{codeblock}{CUDAGridsAndBlocks.cu}
    \begin{lstlisting}[language = C]
#include <stdio.h>
#include <cuda.h>

int main() {
    // Creazione di un blocco di dimensioni 3 x 2 x 1
    dim3 block(3, 2);

    // Creazione di una griglia di dimensioni 4 x 3 x 2
    dim3 grid(4, 3, 2);
}\end{lstlisting}
\end{codeblock}

Il fatto che un vettore specifichi una griglia o un blocco viene definito dall'ordine in cui vengono chiamati quando viene eseguito il kernel. Infatti, si usa la seguente sintassi:

\begin{center}
    \verb|kernel<<<grid, block>>>();|
\end{center}

Nel caso in cui non si determini un valore quando si specifica un vettore di tipo \verb|dim3|, il valore automatico che verrà usato è 1. CUDA tuttavia ritornerà errori nel caso in cui si cerchi di usare un valore illegale per le dimensioni: se ad esempio ci trovassimo su una GPU con compute capability \verb|5.x|, potremmo usare massimo 1024 threads per blocco, e chiedere a CUDA di allocare 2048 threads in un blocco sarebbe illegale.
\nl
Per eseguire una funzione su una GPU NVIDIA, servirà specificare appunto la dimensione dei blocchi e della griglia, e ciò si fa attraverso le seguenti linee di codice:

\begin{codeblock}{CUDAStartKernel.cu}
    \begin{lstlisting}[language = C]
#include <stdio.h>
#include <cuda.h>

__global__ void hello() {
    printf("Hello world");
}

int main() {
    dim3 block(3, 2);
    dim3 grid(4, 3, 2);
    hello <<<grid, block>>>();
    cudaDeviceSynchronize();
    return 0;
}\end{lstlisting}
\end{codeblock}

Analizziamo riga per riga come funziona il codice che abbiamo appena scritto:
\begin{itemize}
    \item [1)] con \verb|#include <cuda.h>| specifichiamo l'inclusione della libreria di CUDA all'interno del nostro programma;
    \item [2)] la direttiva \verb|__global__| determina che la funzione specificata successivamente può essere \textbf{chiamata} sia dall'host (sulla CPU) che sul device esterno (la GPU), tuttavia può essere eseguita solo sul device (quindi sulla GPU);
    \item [3)] la funzione \verb|printf()| è supportata sulle piattaforme CUDA con compute capability superiore a \verb|2.0|, ma è generalmente sconsigliato far stampare elementi in console dal device. In fase di testing è possibile usarlo per fare debugging, ma non andrebbe usato in altri casi;
    \item [4)] con \verb|hello<<<grid, block>>>()| faremmo eseguire la funzione \verb|hello()| su una griglia $4 \times 3 \times 2$ di blocchi di threads la cui dimensione è $3 \times 2 \times 1$;
    \item [5)] la funzione \verb|cudaDeviceSynchronize()| è una sorta di barriera, che fa continuare l'esecuzione del codice solo quando il kernel avrà finito di essere eseguito.
\end{itemize}

Per compilare codice di CUDA è necessario usare \verb|nvcc|, il compilatore di CUDA. Un esempio segue:

\begin{terminal}
    \begin{lstlisting}[style = notexterm]
$ nvcc -arch=2.0 CUDAStartKernel.cu -o CUDAHello
$ ./CUDAHello\end{lstlisting}
\end{terminal}